\documentclass[11pt]{article}
\usepackage{series}

\renewcommand\SeriesName{Modal Triplet Theory}
\renewcommand\SeriesShort{Projection Duality}
\renewcommand\PartRoman{I}
\renewcommand\PartArabic{1}
\renewcommand\PartsInSeries{I}

\title{Wave--Particle Duality as Projection Duality in Modal Triplet Theory\\
\large Pointlike Delta Shadows and Wavelike Spectral Shadows of a Single Coherent Kernel}
\author{Peter Nero}
\date{April 2026}

% Paper-specific notation
\newcommand{\Kadm}{K_{\mathrm{adm}}}
\newcommand{\Badm}{B_{\mathrm{adm}}}
\newcommand{\ta}{\tau_{\mathrm{adm}}}
\newcommand{\epsadm}{\epsilon_{\mathrm{adm}}}
\newcommand{\Leff}{\Lambda_{\mathrm{eff}}}
\newcommand{\Hcoh}{\mathcal H_{\mathrm{coh}}}
\newcommand{\Hinc}{\mathcal H_{\mathrm{inc}}}
\newcommand{\Dab}{D_{ab}}
\newcommand{\Dmat}{D}
\newcommand{\Kdet}{K_{\mathrm{det}}}
\newcommand{\Efx}{E_x}
\newcommand{\calE}{\mathcal E}
\newcommand{\calD}{\mathcal D}
\newcommand{\calA}{\mathcal A}
\newcommand{\calR}{\mathcal R}
\newcommand{\calK}{\mathcal K}
\newcommand{\coh}{\mathrm{coh}}
\newcommand{\inc}{\mathrm{inc}}
\newcommand{\loc}{\mathrm{loc}}
\newcommand{\spec}{\mathrm{spec}}
\newcommand{\MTT}{\mathrm{MTT}}
\newcommand{\diag}{\mathrm{diag}}
\newcommand{\vis}{\mathrm{vis}}
\newcommand{\Meas}{\mathsf M}
\newcommand{\Basin}{\mathcal B}
\newcommand{\CND}{\mathrm{CND}}


\begin{document}
\maketitle
\seriestagline

\begin{abstract}
Wave--particle duality is usually presented as a primitive feature of quantum theory:
microscopic systems propagate as extended wave amplitudes, interfere through coherent
phase relations, and yet appear as localized particle-like events under measurement.
Modal Triplet Theory (MTT) reframes this duality as a projection phenomenon. The
apparent opposition between ``wave'' and ``particle'' is not fundamental. It arises because
one finite coherent-sector excitation admits two downstream shadows: a local delta-like
shadow and a spectral phase-coherent shadow.

In an admissible MTT fixed-point regime, the relevant object is the finite coherent kernel
\[
  \Badm
  =
  P\chi(A)e^{-\ta A}\chi(A)P,
\]
where \(A\) is the fixed-point stabilization operator, \(P\) is the coherent-sector
projector, \(\chi(A)\) is the admissible spectral window, and
\[
  \ta
  =
  \lambda_\ast^{-1}\log(C_Q/\epsadm)
\]
is the damping-selected proper-time/heat-time scale. In the local representation, the kernel
\[
  \Kadm(x,y)=\langle x|\Badm|y\rangle
\]
defines a finite localization profile and recovers the Dirac delta only in the sharp limit.
This is the particle shadow. In the spectral representation,
\[
  \Kadm(x,y)
  =
  \sum_{n\in \coh}
  \chi(\lambda_n)^2 e^{-\ta\lambda_n}
  \phi_n(x)\phi_n^\ast(y),
\]
the same object appears as a weighted coherent modal structure whose phase evolution
produces interference. This is the wave shadow.

The paper strengthens this statement in three ways. First, the delta limit is stated as
a distributional convergence theorem. Second, branch damping is formulated as a Schur
product quantum channel, with the positivity condition \(D\succeq0\) ensuring complete
positivity and trace preservation. Third, finite detector effects are treated as normalized
POVM kernels, so localized detection is a finite survivor-basin selection process rather than
a primitive collapse postulate. The double-slit experiment is then a canonical example:
branch coherence gives interference, while which-way selection damps off-diagonal branch
terms and leaves localized events.

The main result is a projection-duality theorem: wave-like interference and particle-like
localization are not competing ontologies but complementary effective descriptions induced
by different projections of one coherent-sector excitation. Standard quantum mechanics is
recovered in the sharp coherent limit; MTT supplies the deeper projection architecture from
which the dual shadows arise.
\end{abstract}

\tableofcontents

\section{Purpose and claim discipline}

The purpose of this paper is to show that wave--particle duality follows naturally from
the projection architecture of Modal Triplet Theory. The preceding delta-projection sequence
established that Dirac delta distributions should not be read as primitive physical objects
where finite coherent support is relevant. Rather, they arise as singular shadows of finite
admissible kernels. The fixed-point kernel construction then replaced arbitrary smoothing by
a canonical admissible operator
\begin{equation}
  \Badm
  =
  P\chi(A)e^{-\ta A}\chi(A)P.
  \label{eq:badm}
\end{equation}
The present paper applies that same operator to the central interpretive tension of quantum
mechanics: why the same quantum system behaves like an extended wave in interference
experiments and like a localized particle in detection events.

\begin{quote}
\emph{Central claim.}
Wave-like interference and particle-like localization are two downstream projection shadows
of one finite coherent-sector excitation.
\end{quote}

Equivalently,
\[
  \boxed{\text{wave--particle duality}=\text{projection duality}.}
\]

This paper does not claim that standard quantum mechanics is false. It does not deny the
usefulness of wavefunctions, Hilbert spaces, amplitudes, density matrices, POVMs, or
projection-valued idealizations. Instead, it explains why those structures appear: they are
effective descriptions obtained after coherent projection, spectral phase evolution, and
survivor-basin selection.

This paper also does not claim that particles are ordinary extended classical blobs. The MTT
claim is more precise. The underlying object is not a little ball, and it is not a classical wave
spread through space. It is a stable coherent-sector excitation. That excitation can appear
pointlike when probed locally and wavelike when represented spectrally.

Thus
\[
  \text{particle}\neq \text{literal mathematical point},
\]
and
\[
  \text{wave}\neq \text{separate physical substance}.
\]
Rather,
\[
  \text{particle shadow}
  =
  \text{localized delta-limit representation},
\]
while
\[
  \text{wave shadow}
  =
  \text{spectral phase-coherent representation}.
\]
Both are generated by the same admissible MTT kernel.

\begin{remark}[Structural rather than numerical claim]
The paper is a structural synthesis. It does not compute all sector-specific numerical
parameters. Numerical predictions require the concrete sector data \(A\), \(P\), \(\chi\),
\(\lambda_\ast\), \(C_Q\), \(\epsadm\), and the relevant detector or branch-separation metric.
The claim here is that, once those data are supplied by an admissible fixed-point sector, no
independent wave--particle switch or arbitrary smoothing width remains.
\end{remark}

\section{Operational meanings of wave-like and particle-like behavior}

The words ``wave'' and ``particle'' can easily become metaphysical. We therefore fix their
operational meanings.

\begin{definition}[Particle-like behavior]
A coherent excitation is said to display particle-like behavior in a measurement context if:
\begin{enumerate}[label=(\roman*)]
  \item detector outcomes are localized relative to the detector resolution;
  \item repeated trials produce discrete detection records;
  \item the finite detector or source kernel is narrow compared with the measurement scale;
  \item the ideal point-event description is recovered as a delta/PVM limit.
\end{enumerate}
Thus particle-like behavior means localized effective detection, not literal point ontology.
\end{definition}

\begin{definition}[Wave-like behavior]
A coherent excitation is said to display wave-like behavior in an interference context if:
\begin{enumerate}[label=(\roman*)]
  \item the retained coherent sector supports superposition of alternatives or modes;
  \item relative phase information is preserved;
  \item observable probabilities contain off-diagonal cross terms;
  \item these cross terms can be operationally measured as interference or visibility.
\end{enumerate}
Thus wave-like behavior means phase-coherent interference, not a separate classical medium.
\end{definition}

These definitions make the goal precise. We will show that a single admissible coherent
kernel produces both operational behaviors under different downstream projections.

\section{Analytic setting and fixed-point input}

Let \(X\) be a compact smooth Riemannian manifold, or a compact manifold with a
self-adjoint elliptic boundary condition. The scalar notation below may be replaced by
sections of a Hermitian vector bundle. Let
\[
  \calH=L^2(X)
\]
for concreteness.

\begin{remark}[Compact and noncompact charts]
The compact setting is used to state spectral decompositions and kernel convergence cleanly.
Flat noncompact charts, such as \(X=\bbR^d\), are treated by Fourier transform and
tempered-distribution convergence against Schwartz test functions. None of the conceptual
claims below depend on compactness; compactness is a technical convenience for the theorem
statements.
\end{remark}

Let \(A\ge0\) be a nonnegative self-adjoint elliptic or Laplace-type operator with compact
resolvent. Let
\[
  A\phi_n=\lambda_n\phi_n,
  \qquad
  0\le \lambda_0\le\lambda_1\le\cdots,
\]
be an orthonormal spectral resolution, with eigenvalues repeated according to multiplicity.

\begin{assumption}[Admissible fixed-point sector]
We are given:
\begin{enumerate}[label=(\roman*)]
  \item an isolated coherent spectral cluster \(\Omega_{\coh}\subset\sigma(A)\);
  \item the corresponding orthogonal Riesz/spectral projector
  \[
    P=\frac{1}{2\pi i}\oint_{\Gamma}(z-A)^{-1}\,\dd z,
  \]
  where \(\Gamma\) encloses \(\Omega_{\coh}\) and no discarded spectrum;
  \item \(Q=I-P\);
  \item consequently \([P,A]=0\) on the relevant domain and \(P\) preserves the smooth
  spectral domain generated by \(A\);
  \item an incoherent-sector gap \(A|_{\Ran Q}\ge \lambda_\ast I\) for some
  \(\lambda_\ast>0\);
  \item a bounded spectral acceptance window \(0\le\chi\le1\), with \(\chi(A)\) subordinate
  to the same spectral split and commuting with \(P\);
  \item a leakage estimate
  \[
    \norm{Q\Phi_tQ}\le C_Qe^{-\lambda_\ast t},
    \qquad t\ge0,
  \]
  and an admissibility tolerance \(0<\epsadm<1\).
\end{enumerate}
\end{assumption}

The damping-selected proper-time/heat-time scale is
\begin{equation}
  \ta
  =
  \frac{1}{\lambda_\ast}
  \log\left(\frac{C_Q}{\epsadm}\right),
  \label{eq:tauadm}
\end{equation}
whenever \(C_Q>\epsadm\). If \(C_Q\le \epsadm\), this criterion does not force a positive
damping time and one may replace \(\ta\) by
\[
  \ta
  =
  \max\left\{
  0,\,
  \lambda_\ast^{-1}\log(C_Q/\epsadm)
  \right\}.
\]
For the nontrivial case we use \cref{eq:tauadm}.

\begin{definition}[Admissible coherent operator and kernel]
The admissible coherent operator is
\[
  \Badm
  =
  P\chi(A)e^{-\ta A}\chi(A)P.
\]
When \(\Badm\) has a Schwartz kernel, we denote it by
\[
  \Kadm(x,y)=\langle x|\Badm|y\rangle.
\]
\end{definition}

\begin{remark}[Sector identity, not universal identity]
Unless \(P=I\) and \(\chi(A)=I\), \(\Badm\) is not the full identity. It is an identity-like
object on the retained coherent sector, with finite damping and finite spectral acceptance.
This distinction is central: pointlike behavior is a downstream limit, not a primitive
upstream assumption.
\end{remark}


\begin{assumption}[Coherent phase generator]
The wave-like phase evolution is generated by a self-adjoint coherent Hamiltonian
\(H_{\coh}\) on \(\Ran P\), or equivalently by a strongly continuous unitary group
\[
  U_t=e^{-itH_{\coh}}
  \qquad\text{on }\Ran P.
\]
The damping/stabilization operator \(A\) and the coherent phase generator \(H_{\coh}\) need
not be identical. The former selects finite admissible support; the latter transports phases
inside the retained coherent sector.
\end{assumption}

\section{The common-source kernel}

The central mathematical fact is that the same object has local and spectral representations.

\begin{theorem}[Kernel dual-representation theorem]
\label{thm:dual-representation}
Under the admissible fixed-point assumptions, the operator \(\Badm\) has the spectral action
\[
  \Badm f
  =
  \sum_n p_n\chi(\lambda_n)^2e^{-\ta\lambda_n}
  \ip{\phi_n}{f}\phi_n,
\]
where \(p_n\in\{0,1\}\) is the eigenvalue of \(P\) on \(\phi_n\). If \(e^{-\ta A}\) is smoothing
and \(P,\chi(A)\) are spectral functions of \(A\), then \(\Badm\) has a smooth kernel
\[
  \Kadm(x,y)
  =
  \sum_{n:p_n=1}
  \chi(\lambda_n)^2e^{-\ta\lambda_n}
  \phi_n(x)\phi_n^\ast(y).
  \label{eq:spectral-kernel}
\]
Thus the same operator admits:
\begin{enumerate}[label=(\roman*)]
  \item a local kernel representation \(\Kadm(x,y)\);
  \item a spectral modal representation \cref{eq:spectral-kernel}.
\end{enumerate}
\end{theorem}

\begin{proof}
Since \(A\) is self-adjoint with compact resolvent, the spectral theorem gives
\[
  e^{-\ta A}f=\sum_n e^{-\ta\lambda_n}\ip{\phi_n}{f}\phi_n.
\]
Because \(P\) and \(\chi(A)\) commute with \(A\), they are diagonal in the same spectral
resolution in the spectral-projector case. Thus
\[
  P\chi(A)e^{-\ta A}\chi(A)P f
  =
  \sum_n p_n\chi(\lambda_n)^2e^{-\ta\lambda_n}
  \ip{\phi_n}{f}\phi_n.
\]
For \(\ta>0\), \(e^{-\ta A}\) is smoothing for elliptic \(A\) on compact \(X\). Multiplication
by bounded spectral projectors and windows preserves boundedness, and in the purely
spectral-window case preserves the smooth kernel representation on the retained sector.
The displayed kernel is the corresponding Schwartz kernel.
\end{proof}

\begin{remark}[Interpretive separation]
\Cref{thm:dual-representation} is mathematical. The physical interpretation comes later:
the local representation supplies the particle-like localization shadow, while the spectral
representation supplies the wave-like phase-coherence shadow.
\end{remark}

\section{The precise delta limit}

To avoid treating the delta limit as a slogan, we state the sharp-limit theorem explicitly.

Let \(P_N=\mathbf 1_{[0,\Lambda_N]}(A)\), with \(\Lambda_N\to\infty\), and let
\(\tau_N\downarrow0\). Define
\begin{equation}
  B_{N,\tau_N}
  =
  P_N e^{-\tau_N A}P_N.
  \label{eq:BNtau}
\end{equation}

\begin{theorem}[Distributional recovery of the Dirac kernel]
\label{thm:delta-limit}
For every \(f\in C^\infty(X)\),
\[
  B_{N,\tau_N}f\to f
\]
in \(C^\infty(X)\), provided \(\Lambda_N\to\infty\) and \(\tau_N\downarrow0\). Equivalently,
the kernels \(K_{N,\tau_N}(x,y)\) of \(B_{N,\tau_N}\) converge to \(\delta(x-y)\) in the
distributional sense on \(X\times X\).
\end{theorem}

\begin{proof}
Write
\[
  f-B_{N,\tau_N}f
  =
  (I-P_N)f+P_N(I-e^{-\tau_NA})P_Nf.
\]
The first term tends to zero in every Sobolev norm because spectral projectors converge
strongly to the identity on smooth functions and smooth spectral coefficients decay rapidly.
For the second term, fix a Sobolev index \(s\). Choose \(M>s\). Since \(f\in C^\infty(X)\),
\[
  \sum_n (1+\lambda_n)^M\abs{\ip{\phi_n}{f}}^2<\infty.
\]
For each fixed \(n\), \(1-e^{-\tau_N\lambda_n}\to0\). The summand is dominated by a
constant multiple of \((1+\lambda_n)^M\abs{\ip{\phi_n}{f}}^2\) for \(N\) sufficiently large.
Dominated convergence gives convergence to zero in \(H^s\). Since \(s\) is arbitrary,
Sobolev embedding gives convergence in \(C^\infty\). The kernel statement is the standard
distributional form of convergence to the identity operator.
\end{proof}

\begin{corollary}[Particle shadow as sharp local limit]
The ideal point source \(\delta(x-y)\) is not the native finite object in an admissible MTT
sector. It is recovered only when coherent bandwidth becomes complete and the proper-time
width tends to zero.
\end{corollary}

\section{Particle shadow: localized finite kernels}

The particle-like aspect appears when the coherent kernel is probed locally. An ideal point
source centered at \(x_0\) is written
\[
  \rho_{x_0}^{\mathrm{point}}(x)=\delta(x-x_0).
\]
MTT replaces the point source by the finite local response amplitude
\[
  k_{x_0}^{\mathrm{MTT}}(x)
  =
  \Kadm(x,x_0).
\]
For positive heat kernels this may itself be read as a smeared source density. In general,
however, \(\Kadm(x,x_0)\) is an amplitude or response kernel and need not be nonnegative.
Detection probabilities are obtained from associated effects such as
\[
  E_{x_0}=|k_{x_0}\rangle\langle k_{x_0}|,
\]
not by assuming that the kernel value itself is always a probability density.

\begin{definition}[Effective pointlikeness]
A finite localization kernel is effectively pointlike at detector resolution \(\ell_{\det}\) if its
characteristic width \(\ell_{\coh}\) satisfies
\[
  \ell_{\coh}\ll \ell_{\det}.
\]
In flat heat-kernel models, \(\ell_{\coh}\sim 2\sqrt{\ta}\), so effective pointlikeness means
\[
  2\sqrt{\ta}\ll \ell_{\det}.
\]
Equivalently, if \(\Leff=\ta^{-1/2}\), pointlike behavior at probe energy \(E\) requires
\[
  E\ll \Leff,
\]
up to experiment-dependent tolerance factors.
\end{definition}

\begin{example}[Flat heat-kernel profile]
On \(X=\bbR^d\) with \(A=-\Delta\), \(P=I\), and \(\chi=1\), the heat kernel is
\[
  K_\tau(x,y)
  =
  (4\pi\tau)^{-d/2}
  \exp\left(-\frac{\abs{x-y}^2}{4\tau}\right).
\]
As \(\tau\downarrow0\), \(K_\tau(x,y)\to\delta(x-y)\) distributionally. For finite \(\tau\),
the object has finite coherent width. It appears pointlike only to probes that cannot resolve
that width.
\end{example}

The particle is therefore not a literal mathematical point. It is the local delta shadow of a
finite coherent excitation.

\section{Wave shadow: spectral phase coherence}

The wave-like aspect appears when the same coherent excitation is represented spectrally.
A retained coherent state may be written
\[
  \psi_{\coh}(x)
  =
  \sum_{n\in\coh}c_n\phi_n(x).
\]
If \(H_{\coh}\) is the coherent-sector phase generator, then
\[
  \psi_{\coh}(t)
  =
  e^{-itH_{\coh}}\psi_{\coh}(0),
\]
or, in modes,
\[
  \psi_{\coh}(x,t)
  =
  \sum_{n\in\coh}
  c_ne^{-i\omega_nt}\phi_n(x).
\]

\begin{proposition}[Interference from retained phase relations]
Let a coherent excitation decompose into two retained branches,
\[
  \psi=\psi_a+\psi_b.
\]
If both branches remain in the same coherent sector and their relative phase is retained, then
the intensity contains the cross term
\[
  \abs{\psi}^2
  =
  \abs{\psi_a}^2+\abs{\psi_b}^2
  +
  2\operatorname{Re}(\psi_a^\ast\psi_b).
\]
Thus wave-like interference is the operational signature of retained off-diagonal coherence.
\end{proposition}

\begin{proof}
The identity follows by expanding \((\psi_a+\psi_b)^\ast(\psi_a+\psi_b)\). The physical
content is that the cross term is observable only when the downstream description retains
the relative phase between the two branches.
\end{proof}

The wave is therefore not a separate substance. It is the spectral phase shadow of the same
coherent excitation whose local shadow appears particle-like.

\section{Branch damping as a quantum channel}

The informal rule
\[
  \rho_{ab}\mapsto D_{ab}\rho_{ab}
\]
must be mathematically valid. In particular, it should define a completely positive
trace-preserving map on the retained branch density matrix.

Let \(\{a\}_{a=1}^N\) label a finite branch decomposition and let
\[
  \rho=(\rho_{ab})_{a,b=1}^N
\]
be the density matrix in that branch basis.

\begin{definition}[Schur branch-damping map]
Let \(D=(D_{ab})\) be a complex \(N\times N\) matrix. Define
\[
  \calE_D(\rho)=D\circ \rho,
\]
where \(\circ\) denotes entrywise, or Schur, product:
\[
  (\calE_D(\rho))_{ab}=D_{ab}\rho_{ab}.
\]
\end{definition}

\begin{theorem}[Complete positivity of admissible branch damping]
\label{thm:schur-channel}
If \(D\succeq0\) is positive semidefinite and \(D_{aa}=1\) for all \(a\), then
\[
  \calE_D(\rho)=D\circ\rho
\]
is completely positive and trace preserving.
\end{theorem}

\begin{proof}
This is the standard Schur-product construction of a dephasing channel; it uses the Schur
product theorem and the usual Kraus characterization of completely positive maps
\cite{HornJohnson,Paulsen,NielsenChuang}. Since \(D\succeq0\), there exist vectors
\(v_a\) in a finite-dimensional Hilbert space such that \(D_{ab}=\ip{v_b}{v_a}\). Let
\(v_a^\mu\) be their components in an orthonormal basis
and define diagonal Kraus operators
\[
  M_\mu=\sum_a v_a^\mu |a\rangle\langle a|.
\]
Then
\[
  \sum_\mu M_\mu \rho M_\mu^\ast
\]
has matrix entries
\[
  \sum_\mu v_a^\mu \rho_{ab}\overline{v_b^\mu}
  =
  \ip{v_b}{v_a}\rho_{ab}
  =
  D_{ab}\rho_{ab}.
\]
Thus \(\calE_D\) is completely positive. Since \(D_{aa}=1\), the diagonal entries are unchanged,
so \(\tr(\calE_D(\rho))=\tr(\rho)\).
\end{proof}

\begin{corollary}[Consistency condition for MTT damping factors]
A proposed MTT branch-damping family \(D_{ab}\) is an admissible retained-sector quantum
operation if
\[
  D\succeq0,
  \qquad
  D_{aa}=1.
\]
Thus damping factors are not arbitrary visibility multipliers. They must define a positive
Schur multiplier on branch density matrices.
\end{corollary}

\begin{proposition}[Gaussian branch-distance damping is admissible]
\label{prop:gaussian-damping}
Let \(r_a\) be points in a real Hilbert space and let
\[
  D_{ab}
  =
  \exp\left(-\frac{\tau}{2}\norm{r_a-r_b}^2\right),
  \qquad \tau\ge0.
\]
Then \(D\succeq0\) and \(D_{aa}=1\). Hence \(D\circ\rho\) is a completely positive
trace-preserving branch-damping map.
\end{proposition}

\begin{proof}
The Gaussian kernel \(G_\tau(r,s)=\exp[-\tau\norm{r-s}^2/2]\) is positive definite on real
Hilbert spaces. For finite-dimensional branch data this follows by the Fourier representation
of the Gaussian after embedding the finite span of the \(r_a\) into \(\bbR^m\). Therefore the
finite Gram matrix \(D_{ab}=G_\tau(r_a,r_b)\) is positive semidefinite. The diagonal entries
are \(D_{aa}=1\).
\end{proof}

\begin{definition}[Conditionally negative branch separation]
A symmetric matrix \(\Lambda=(\Lambda_{ab})\) with \(\Lambda_{aa}=0\) is conditionally
negative definite if
\[
  \sum_{a,b}\overline{c_a}c_b\Lambda_{ab}\le0
  \qquad
  \text{whenever }\sum_a c_a=0.
\]
\end{definition}

\begin{proposition}[Schoenberg admissibility criterion]
\label{prop:schoenberg}
If \(\Lambda\) is conditionally negative definite, then for every \(\tau\ge0\)
\[
  D_{ab}=e^{-\tau\Lambda_{ab}}
\]
is positive semidefinite and satisfies \(D_{aa}=1\). Hence \(D\circ\rho\) is an admissible
branch-damping channel.
\end{proposition}

\begin{proof}
This is Schoenberg's theorem on conditionally negative kernels and positive definite
exponentials. The Gaussian distance kernel in \cref{prop:gaussian-damping} is the special
case \(\Lambda_{ab}=\norm{r_a-r_b}^2/2\).
\end{proof}

\begin{remark}[Branch-separation scales]
A scalar expression \(D_{ab}=e^{-\ta\Lambda_{ab}}\) is safe when \(\Lambda\) is conditionally
negative definite, or when it arises from a positive kernel such as \cref{prop:gaussian-damping}.
In general, the matrix \(D\), not merely the pairwise numbers \(\Lambda_{ab}\), must satisfy
complete-positivity constraints.
\end{remark}

\section{Measurement as finite survivor-basin selection}

In ordinary notation, an ideal position measurement is represented by the sharp projector
\[
  |x\rangle\langle x|.
\]
In MTT this is a singular idealization. A physical detector has finite resolution, finite
disturbance strength, finite coherence capacity, and finite admissible support. Its effects are
therefore represented by a POVM, in the standard operational sense of generalized quantum
measurement \cite{NielsenChuang}.

\begin{definition}[Finite detector-effect kernel]
Let \(k_x\in\Hcoh\) be a detector response vector centered at outcome \(x\). Define
\[
  E_x=|k_x\rangle\langle k_x|.
\]
In coordinate notation,
\[
  E_x(y,z)=K_{\det}(y,x)K_{\det}^\ast(z,x).
\]
The detector kernel is a detector-sector instance of the same admissible-kernel construction:
\[
  B_{\det}
  =
  P_{\det}\chi_{\det}(A_{\det})
  e^{-\tau_{\det}A_{\det}}
  \chi_{\det}(A_{\det})P_{\det}.
\]
Thus \(K_{\det}\) is not a new arbitrary smearing profile; it is the measurement-device
realization of the finite coherent kernel.
\end{definition}

\begin{assumption}[Retained-sector POVM normalization]
For a closed retained detector model,
\[
  \int E_x\,\dd x=I_{\coh}.
  \label{eq:povm-normalization}
\]
For an open detector model with unobserved discarded channels, one has instead
\[
  \int E_x\,\dd x\le I_{\coh}.
\]
The closed case is used for probability-conserving detection records.
\end{assumption}

\begin{proposition}[Probability conservation for finite detector effects]
If \(\rho\) is a retained-sector density matrix and \(\{E_x\}\) satisfies
\[
  \int E_x\,\dd x=I_{\coh},
\]
then
\[
  p(x)=\tr(\rho E_x)
\]
is normalized:
\[
  \int p(x)\,\dd x=1.
\]
\end{proposition}

\begin{proof}
Using linearity and the POVM normalization,
\[
  \int p(x)\,\dd x
  =
  \int \tr(\rho E_x)\,\dd x
  =
  \tr\left(\rho\int E_x\,\dd x\right)
  =
  \tr(\rho I_{\coh})
  =
  \tr(\rho)
  =
  1.
\]
\end{proof}

Measurement is therefore not a primitive jump from wave to particle. It is the sequence
\[
  \text{coherent evolution}
  \to
  \text{localized disturbance}
  \to
  \text{projection}
  \to
  \text{survivor-basin stabilization}.
\]
The sharp PVM is recovered only in the ideal limit \(E_x\to |x\rangle\langle x|\).

\section{Device-dependent disturbance and exact records}

Different measurement devices do not merely reveal the same downstream shadow with different
labels. They implement different finite couplings, different detector kernels, and sometimes
different survivor-basin partitions. In MTT a measurement arrangement \(\Meas\) is specified by
device-sector data
\[
  (A_{\Meas},P_{\Meas},\chi_{\Meas},\tau_{\Meas},\{E_i^{(\Meas)}\}_{i\in I}).
\]
It therefore determines both a POVM and a branch-distance kernel. For branches \(a,b\), the
same upstream coherent excitation may acquire different damping factors under different
apparatuses:
\[
  D_{ab}^{(\Meas)}
  =
  \exp[-\tau_{\Meas}\Lambda_{ab}^{(\Meas)}],
\]
subject to the complete-positivity condition of \cref{thm:schur-channel} or
\cref{prop:schoenberg}.

\begin{definition}[Measurement context]
A measurement context is an admissible device-sector projection
\[
  \Meas:\rho\longmapsto
  \{p_i=\tr(\rho E_i^{(\Meas)})\}_{i\in I}
\]
together with a stabilization map assigning the post-interaction configuration to survivor
basins \(\{B_i^{(\Meas)}\}\).
\end{definition}

\begin{proposition}[Device dependence of particle-like records]
Let two devices \(\Meas\) and \(\Meas'\) couple to the same coherent input state but induce
different detector effects or different branch-distance kernels. Then they need not preserve
the same interference visibility, nor select the same effective branch basis. What is exact is
the stabilized record inside the chosen measurement context:
\[
  \Psi\longmapsto B_i^{(\Meas)}
\]
after disturbance, projection, and re-coherence.
\end{proposition}

\begin{proof}
The probabilities and coherences measured by a device are determined by its POVM effects
and branch-damping matrix. If \(E_i^{(\Meas)}\neq E_i^{(\Meas')}\) or
\(D^{(\Meas)}\neq D^{(\Meas')}\), the downstream statistics and visibility differ. Once a
particular device coupling drives the configuration into a survivor basin, the stabilized
record is discrete within that context. Thus exact records are contextual stabilized outcomes,
not non-disturbing revelations of a context-independent point property.
\end{proof}

\begin{remark}[Weak, strong, and quantum-eraser regimes]
A weak which-way device has \(D_{12}^{(\Meas)}\) close to one and only partially reduces
visibility. A strong which-way device has \(D_{12}^{(\Meas)}\) close to zero and yields
particle-like records. A quantum-eraser arrangement changes the final device projection so
that the branch marker is measured in a basis where conditional coherences can reappear.
In MTT language, the eraser does not undo a literal past trajectory; it changes the downstream
survivor-basin partition used to read the branch-marked state.
\end{remark}



\section{Double-slit experiment in operator form}

The double-slit experiment is the cleanest illustration of projection duality.

Let \(|\psi_1\rangle\) and \(|\psi_2\rangle\) be the two branch amplitudes after the slits,
including their complex weights. The normalized outgoing state is
\[
  |\psi\rangle
  =
  \frac{|\psi_1\rangle+|\psi_2\rangle}{\sqrt{\mathcal N}},
  \qquad
  \mathcal N=\norm{\psi_1+\psi_2}^2,
\]
and, to avoid clutter, the normalization factor is absorbed into the branch amplitudes below. The pre-detection density matrix decomposes into
\[
  \rho
  =
  \rho_{11}+\rho_{22}+\rho_{12}+\rho_{21},
\]
where
\[
  \rho_{ab}=|\psi_a\rangle\langle\psi_b|.
\]

A branch-damping channel maps
\[
  \rho
  \mapsto
  \rho'
  =
  \rho_{11}+\rho_{22}
  +
  D_{12}\rho_{12}
  +
  D_{21}\rho_{21},
\]
with \(D_{11}=D_{22}=1\), \(D_{21}=D_{12}^\ast\), and \(D\succeq0\).

For a screen effect \(E_x\), the detected intensity is
\[
  I_{\MTT}(x)
  =
  \tr(\rho' E_x).
\]
In the ideal narrow-screen limit this gives
\begin{equation}
  I_{\MTT}(x)
  =
  \abs{\psi_1(x)}^2
  +
  \abs{\psi_2(x)}^2
  +
  2D_{12}\operatorname{Re}\left[\psi_1^\ast(x)\psi_2(x)\right],
  \label{eq:slit-intensity}
\end{equation}
for real \(D_{12}\in[0,1]\). More generally, the phase of \(D_{12}\) shifts the interference
phase.

\begin{proof}[Derivation of \cref{eq:slit-intensity}]
Using \(E_x\approx |x\rangle\langle x|\),
\[
  \tr(\rho_{11}E_x)=\abs{\psi_1(x)}^2,\qquad
  \tr(\rho_{22}E_x)=\abs{\psi_2(x)}^2,
\]
and
\[
  \tr(\rho_{12}E_x)=\psi_2^\ast(x)\psi_1(x),
  \qquad
  \tr(\rho_{21}E_x)=\psi_1^\ast(x)\psi_2(x).
\]
Adding the damped off-diagonal terms gives \cref{eq:slit-intensity}.
\end{proof}

No which-way selection corresponds to
\[
  D_{12}\approx1,
\]
so the interference term survives. Strong which-way selection corresponds to
\[
  D_{12}\approx0,
\]
so the interference term is suppressed and the downstream record becomes particle-like.

\subsection{Visibility}

For symmetric slits with equal single-branch intensities and ideal visibility \(V_0\), define
\[
  V=\frac{I_{\max}-I_{\min}}{I_{\max}+I_{\min}}.
\]

\begin{proposition}[Visibility damping]
\label{prop:visibility}
In the symmetric two-branch case with real \(0\le D_{12}\le1\),
\[
  V_{\MTT}=D_{12}V_0.
\]
\end{proposition}

\begin{proof}
In a symmetric two-branch interference pattern the undamped intensity may be written
locally as
\[
  I_0(x)=\bar I(x)\bigl[1+V_0\cos\varphi(x)\bigr],
\]
where \(\bar I(x)\) is the slowly varying envelope and \(V_0\) is the ideal fringe
visibility. The Schur damping channel multiplies only the off-diagonal coherence and hence
only the modulation term. Therefore
\[
  I_{\MTT}(x)=\bar I(x)\bigl[1+D_{12}V_0\cos\varphi(x)\bigr].
\]
Consequently
\[
  \frac{I_{\max}-I_{\min}}{I_{\max}+I_{\min}}=D_{12}V_0.
\]
\end{proof}

This gives an operational bridge:
\[
  D_{12}\approx1
  \Rightarrow
  \text{wave-like interference},
\]
while
\[
  D_{12}\approx0
  \Rightarrow
  \text{particle-like which-way behavior}.
\]
In standard two-path language this is the same tradeoff measured by visibility and which-way
distinguishability; the MTT contribution is to identify \(D_{12}\) with an admissible
branch-coherence survival factor rather than with an independent postulate \cite{Englert}.

\begin{remark}[Not ordinary propagation decoherence]
The factor \(D_{12}\) in this paper is a retained-sector branch-coherence factor. It is not
assumed to be a propagation-length decoherence exponent. Propagation-dependent damping
would require an additional open-system or environmental dynamics not assumed here.
\end{remark}

\section{Further experimental modules}

The double-slit experiment is the minimal example. The same projection-duality pattern also
organizes several standard quantum experiments.

\subsection{Stern--Gerlach splitting}

A Stern--Gerlach device does not test an abstract spin property without context. It supplies
a magnetic-field gradient, hence a measurement context \(\Meas_{\hat n}\) associated with an
axis \(\hat n\). The branch basis is
\[
  \{|\uparrow_{\hat n}\rangle,|\downarrow_{\hat n}\rangle\},
\]
and the spatial detector separates the two spin-correlated beams. In MTT terms the apparatus
creates two survivor basins,
\[
  B_{\uparrow}^{(\hat n)},\qquad B_{\downarrow}^{(\hat n)},
\]
and a strong device drives
\[
  D_{\uparrow\downarrow}^{(\hat n)}\approx0.
\]
The outcome is exact within the chosen apparatus context: a spot in the upper or lower basin.
Changing the magnet orientation changes the measurement context and therefore the branch
basis selected by the device. Thus exact records are compatible with contextual disturbance.

\subsection{Mach--Zehnder interferometer}

A Mach--Zehnder interferometer is a controlled two-branch system. With no which-path marker,
the two path branches remain coherent and recombine:
\[
  D_{12}\approx1.
\]
A phase shifter changes the relative phase \(\varphi\), producing output probabilities of the
form
\[
  p_{\pm}=\frac12\bigl(1\pm V_0\cos\varphi\bigr)
\]
in the ideal balanced case. A which-path marker changes the branch-damping matrix, yielding
\[
  p_{\pm}^{\MTT}
  =
  \frac12\bigl(1\pm D_{12}V_0\cos\varphi\bigr).
\]
The apparatus therefore interpolates continuously between wave-like recombination and
particle-like path records.

\subsection{Delayed-choice and quantum eraser arrangements}

Delayed-choice experiments change the final projection context after the branches have been
prepared. In MTT this does not require retrocausal alteration of a past particle path. The
branch-marked state retains or loses operational interference depending on the final
survivor-basin partition. A quantum eraser measures the marker in a basis that groups branch
records into coherent conditional subensembles. Thus unconditional interference may remain
absent while conditional fringes reappear. The MTT distinction is:
\[
  \text{unconditional record}
  \neq
  \text{conditional survivor-basin partition}.
\]

\subsection{Matter-wave interferometry and diffraction}

Electron, neutron, atom, and molecule diffraction are direct examples of the same structure.
The local detection event is particle-like because the detector effect is narrow, while the
propagating object is wave-like because retained modal phases remain coherent across paths or
grating openings. Finite mass changes the coherent phase generator \(H_{\coh}\), not the
projection-duality architecture:
\[
  \psi_{\coh}(t)=e^{-itH_{\coh}}\psi_{\coh}(0).
\]

\subsection{Hong--Ou--Mandel-type indistinguishability}

Two-particle interference experiments emphasize that interference is not merely a property of
single-particle trajectories. The relevant branches may be alternatives in a multi-particle
configuration space. If the alternatives are indistinguishable within the detector context, their
off-diagonal coherences survive. If auxiliary labels make them distinguishable, the corresponding
branch-damping factor decreases. Thus the same Schur-channel logic applies to configuration
branches, not only to spatial paths.

\section{The wave--particle projection-duality theorem}

We can now state the central result.

\begin{theorem}[Wave--particle projection duality]
\label{thm:projection-duality}
Let \((A,P,\chi,\ta)\) define an admissible fixed-point coherent sector with operator
\[
  \Badm=P\chi(A)e^{-\ta A}\chi(A)P.
\]
Assume:
\begin{enumerate}[label=(\roman*)]
  \item \(A\) is nonnegative self-adjoint and has a coherent spectral sector selected by \(P\);
  \item \(\chi(A)\) is subordinate to the admissible coherent spectral window;
  \item \(\ta\) is fixed by discarded-sector damping and admissibility tolerance;
  \item coherent evolution preserves relative phase inside \(\Ran P\);
  \item detector effects form finite POVMs on the retained sector;
  \item branch damping is implemented by an admissible Schur channel \(D\circ\rho\).
\end{enumerate}
Then the same coherent-sector excitation has two downstream operational shadows:
\begin{enumerate}[label=(\alph*)]
  \item a particle-like shadow, given by local finite localization kernels that recover
  \(\delta(x-y)\) in the sharp limit;
  \item a wave-like shadow, given by spectral coherent modes whose retained relative phases
  generate interference.
\end{enumerate}
Measurement or which-way selection damps off-diagonal branch coherences. Hence
wave-like interference and particle-like localization are not distinct ontological primitives.
They are complementary projection shadows of one finite coherent-sector excitation.
\end{theorem}

\begin{proof}
The local statement follows from \cref{thm:delta-limit}: finite kernels converge to the
Dirac identity only in the joint sharp limit. Therefore local probing of a narrow finite kernel
yields effective particle-like localization.

The spectral statement follows from \cref{thm:dual-representation}: the same operator
\(\Badm\) has a retained modal expansion. Coherent evolution multiplies retained modes by
phases \(e^{-i\omega_nt}\), and superpositions of retained branches produce cross terms in
observable probabilities.

The measurement statement follows from the POVM normalization result and
\cref{thm:schur-channel}. Finite detector effects produce normalized probabilities, while
admissible Schur damping suppresses off-diagonal branch coherences without violating
complete positivity or trace preservation.

Thus the same coherent-sector excitation yields wave-like or particle-like behavior depending
on which downstream projection context is applied. This proves the projection-duality claim.
\end{proof}

\[
  \boxed{\text{wave--particle duality}=\text{projection duality}.}
\]

\section{Recovery of standard quantum mechanics}

MTT does not discard standard quantum mechanics. It recovers it in appropriate limits.

\begin{proposition}[Recovery limits]
The standard formulas are recovered under the following limiting conditions:
\begin{enumerate}[label=(\roman*)]
  \item If \(D_{ab}=1\) for all retained branches, interference is fully coherent.
  \item If \(D_{ab}=0\) for \(a\ne b\), the retained state reduces to an incoherent branch
  mixture.
  \item If \(K_{\det}\) tends to a delta detector response, \(E_x\to |x\rangle\langle x|\).
  \item If \(P\chi(A)^2P\to I\) and \(\ta\to0\), then \(\Kadm(x,y)\to\delta(x-y)\).
\end{enumerate}
\end{proposition}

Thus the ordinary wavefunction appears as the coherent spectral shadow,
\[
  \psi(x,t)=\sum_n c_ne^{-iE_nt}\phi_n(x),
\]
the ordinary point measurement appears as the sharp detector limit,
\[
  E_x\to |x\rangle\langle x|,
\]
and the ordinary point source appears as the sharp kernel limit,
\[
  \Kadm(x,y)\to\delta(x-y).
\]

\section{Compatibility with basin-measure and measurement papers}

This paper is designed to sit between two earlier MTT measurement results. The first
identifies Born probabilities with basin weights of an invariant basin-measure functional.
The second interprets measurement as localized disturbance followed by re-coherence,
projection, and stabilization into one admissible basin. The present paper adds a different
piece: it identifies the wave--particle contrast with the survival or suppression of
off-diagonal branch coherence.

\subsection{Outcome frequencies versus interference visibility}

The branch-damping channel used above controls interference visibility. It does not by
itself assign outcome frequencies. Outcome frequencies are supplied by the basin-measure
functional. Let \(\{B_i\}_{i\in I}\) be the stabilized outcome basins inside an admissible
coherent slab, and let \(\mu\) be the invariant basin measure. Then the basin-measure
probability of outcome \(i\) is
\[
  p_i
  =
  \frac{\mu(B_i)}{\sum_{j\in I}\mu(B_j)}.
\]
The basin-measure statement requires the usual measurement assumptions: a stationary
slab or ensemble measure \(\mu\), a measurable partition by survivor basins
\(\{B_i\}_{i\in I}\), and almost-sure finite capture into exactly one stabilized basin. Under
those hypotheses, \(p_i\) is a basin-measure weight rather than an interference-visibility
factor.

The role of \(D_{ab}\) is different: it controls whether coherent alternatives remain
visible as interference before stabilization. Thus the duality paper separates two questions
that are often conflated:
\[
\begin{aligned}
  &\text{which stabilized outcome occurs?}\\
  &\text{how much branch coherence remains before selection?}
\end{aligned}
\]

\begin{proposition}[Compatibility with basin-measure probabilities]
Assume a normalized detector POVM \(\{E_i\}_{i\in I}\) whose effects represent finite
survivor basins. Assume also a stationary slab measure \(\mu\), a measurable survivor-basin
partition \(\{B_i\}_{i\in I}\), and almost-sure finite capture into one basin, so that
\[
  p_i=\mu(B_i)/\sum_j\mu(B_j).
\]
If the pre-selection coherence between alternatives is damped by a CPTP Schur channel
\(\calE_D\), then branch visibility is modified by \(D\), while the stabilized outcome
weights are still computed by the basin measures of the survivor basins.
\end{proposition}

\begin{proof}
The Schur channel acts on off-diagonal branch coherences and preserves trace by
\cref{thm:schur-channel}. The basin-measure rule assigns probabilities after the
stabilizing basin-capture event and depends only on the measurable basin partition.
Therefore Schur damping changes interference terms in pre-selection observables but does
not replace the basin-measure rule for stabilized outcomes.
\end{proof}

\subsection{Disturbance strength and smooth wave--particle crossover}

The disturbance-stabilization paper describes double-slit behavior as a smooth crossover:
weak which-way disturbance leaves the phase-sector floor small and interference survives;
strong disturbance pushes the configuration across an admissibility threshold and stabilizes
single-path basins. The present formalism expresses the same statement by letting the
branch-separation scale depend on the detector disturbance strength \(\delta\):
\[
  D_{12}(\delta)
  =
  \exp\bigl[-\tau_{\mathrm{det}}\Lambda_{12}^{\mathrm{det}}(\delta)\bigr].
\]
If \(\Lambda_{12}^{\mathrm{det}}(\delta)\) is nondecreasing, then \(D_{12}(\delta)\) is
nonincreasing. Fringe visibility therefore decreases continuously as which-way
distinguishability increases:
\[
  V_{\MTT}(\delta)
  =
  D_{12}(\delta)V_0.
\]
This is the rigorous version of the earlier qualitative claim that wave--particle behavior is
not an ontological switch but a controlled transition between admissible projection regimes.

\begin{remark}[OU floors and finite detector width]
The disturbance-stabilization analysis also implies that detector kernels generally possess
irreducible disturbance--damping floors. In a linearized OU regime one expects a lower-width
estimate of the form
\[
  \ell_{\mathrm{eff}}^2\gtrsim \frac{\delta}{2\gamma},
\]
where \(\delta\) is disturbance power and \(\gamma\) is the damping margin. Thus the sharp
PVM limit is useful notation, but finite POVMs are the physically native measurement
objects in MTT.
\end{remark}


\section{Relation to existing accounts}

The claim is not that no previous framework addressed wave--particle duality. Many did.

\subsection{Complementarity}

Standard complementarity, beginning with Bohr's analysis \cite{Bohr}, states that wave-like and particle-like descriptions are revealed by
different experimental arrangements. MTT preserves this insight but gives it a projection
mechanism:
\[
  \Psi_{\coh}
  \to
  \begin{cases}
  \text{spectral phase shadow},\\
  \text{localized survivor-basin shadow}.
  \end{cases}
\]
The descriptions are complementary because downstream projections are non-invertible.

\subsection{Decoherence}

Decoherence theory \cite{Zurek,Schlosshauer} explains suppression of off-diagonal density-matrix terms through environmental
entanglement. MTT uses a similar mathematical form,
\[
  \rho_{ab}\mapsto D_{ab}\rho_{ab},
\]
but interprets admissible damping as a projection/selection effect tied to coherent-sector
data. Environmental decoherence may be one physical mechanism realizing such damping,
but it is not the only admissibility mechanism.

\subsection{Pilot-wave theory}

Pilot-wave theory \cite{Bohm} keeps both a particle configuration and a guiding wave as real structures.
MTT instead treats local particle behavior and spectral wave behavior as two downstream
shadows of one coherent-sector excitation. The ontology is therefore not a particle-plus-wave
dual ontology, but a projection-duality account.

\subsection{Quantum field theory}

Quantum field theory already treats particles as excitations of fields and uses local
distributions, propagators, and detector models. MTT is compatible with that operational
success, but reinterprets point-local distributions as sharp limits of finite coherent kernels.

\subsection{String and extended-object accounts}

String theory \cite{Polchinski} replaces point particles with extended strings, whose low-resolution behavior
can appear pointlike and whose vibrational modes define spectra. MTT is different: the common
source is not a fundamental one-dimensional object but a finite admissible coherent kernel
selected by projection and damping data. The pointlike and wavelike aspects are kernel
shadows, not merely geometric extension versus vibration.

\begin{remark}[Relation to earlier MTT measurement papers]
Earlier MTT measurement papers supply the basin-measure and disturbance-stabilization
sides of the story: probabilities arise from relative basin weights, finite POVMs are native
survivor-basin effects, and ideal collapse is a singular sharp-basin limit. The present paper
does not rederive those results. It uses them to explain why the same coherent excitation
has a wave-like phase-coherence shadow before selection and a particle-like localized
survivor shadow after selection.
\end{remark}

\begin{remark}[Novelty statement]
The novelty is not that finite kernels, wavefunctions, or decoherence factors exist. The
novelty is the unified selection architecture: the same admissible coherent kernel has a
local delta-shadow, a spectral phase-shadow, and a measurement-selection shadow, with its
width fixed by damping/admissibility data rather than by an external regulator choice.
\end{remark}

\section{Relation to propagators and point interactions}

The same structure appears in propagators. In a flat Euclidean coherent chart with
\(A=-\Delta\), the MTT-corrected scalar propagator is
\[
  \Delta_{\mathrm{adm}}(k)
  =
  \frac{e^{-\ta k^2}}{k^2+m^2}.
\]
This is the propagator-level expression of finite coherent support.

Point sources are replaced by coherent kernels:
\[
  \delta(x-y)\rightsquigarrow \Kadm(x,y).
\]
Point contact interactions are similarly replaced by finite coherent-overlap vertices.
A local \(n\)-point contact profile
\[
  \int_X \prod_{j=1}^n \delta(x-x_j)\,\dd x
\]
is replaced by
\[
  \int_X \prod_{j=1}^n \Kadm(x,x_j)\,\dd x.
\]
Thus the point-particle limit, local source limit, and contact-interaction limit are singular
shadows of finite coherent support.

\section{Relation to branch coherence in other sectors}

The double-slit branch damping structure is not special to spatial slits. It is a general
coherent-branch phenomenon.

For neutrino mass branches, the analogous admissible damping factor is
\[
  D_{ij}^{\nu}
  =
  \exp(-\tau_{\nu,\mathrm{adm}}\abs{\Delta m_{ij}^2}),
\]
or, after inserting the neutrino-sector damping scale,
\[
  D_{ij}^{\nu}
  =
  \exp\left[
  -
  \frac{\abs{\Delta m_{ij}^2}}{\lambda_{\ast,\nu}}
  \log\left(\frac{C_{Q,\nu}L_\nu}{d_{\nu,\mathrm{basin}}}\right)
  \right].
\]
For double-slit spatial branches, the analogous expression is
\[
  D_{12}^{\mathrm{slit}}
  =
  \exp(-\tau_{\mathrm{det}}\Lambda_{12}^{\mathrm{det}}),
\]
provided the resulting branch matrix is positive semidefinite or arises from a positive
branch-distance kernel. The operators and scales differ, but the structure is the same:
\[
\begin{aligned}
  \text{coherent branch superposition}
  &\longrightarrow
  \text{phase evolution} \\
  &\longrightarrow
  \text{possible damping by admissibility or selection}.
\end{aligned}
\]

\section{Failure modes and domain of validity}

Stating the failure modes is part of the rigor of the result.

\begin{enumerate}[label=(\roman*)]
  \item If no isolated coherent spectral sector exists, the projector \(P\) is not canonically
  defined.
  \item If \(P\) is not a spectral/Riesz projector or otherwise fails to preserve the smooth
  domain of \(A\), smooth-kernel claims require additional regularity assumptions.
  \item If the discarded-sector gap vanishes, \(\lambda_\ast=0\), the damping time may diverge
  and no finite admissible kernel is selected by this mechanism.
  \item If \(\chi(A)\) is tuned independently of the spectral split, it becomes a regulator choice
  rather than an admissible MTT window.
  \item If \(H_{\coh}\) does not generate coherent unitary phase evolution on \(\Ran P\), the
  wave-shadow portion must be replaced by the appropriate open-system evolution.
  \item If detector effects do not form a normalized POVM, probability either leaks into
  unobserved channels or must be renormalized by an explicit conditioning rule.
  \item If a proposed branch-damping matrix \(D\) is not positive semidefinite, the map
  \(\rho\mapsto D\circ\rho\) is not a completely positive quantum operation.
  \item If the finite kernel width is experimentally resolvable, the point-particle approximation
  fails and finite-width corrections should appear.
  \item If the measurement device is changed, the relevant POVM, branch basis, and damping
  matrix may change; exact records are exact only inside the specified device context.
\end{enumerate}

These are not defects of the theorem. They are its boundary conditions.

\section{Phenomenological consequences}

The paper is primarily structural, but it gives immediate scaling consequences.

\subsection{Pointlike bounds}

If no substructure is resolved up to an energy scale \(E\), then the effective coherence scale
must satisfy
\[
  \Leff\gtrsim E.
\]
More carefully, if deviations are bounded at fractional level \(\eta\), then in flat sectors
\[
  \ta E^2\lesssim \eta.
\]

\subsection{Interference visibility}

For a two-branch interference experiment,
\[
  V_{\MTT}=D_{ab}V_0.
\]
Thus visibility constraints bound the product
\[
  \tau_{\mathrm{sel}}\Lambda_{ab}
\]
or, in distance-kernel realizations, the distinguishability distance between branches.

\subsection{Detector resolution}

Pointlike detection requires
\[
  2\sqrt{\tau_{\mathrm{det}}}\ll \ell_{\mathrm{det}}
\]
in heat-kernel models. If the detector kernel becomes comparable to the detector resolution,
finite-width corrections to ideal point events may become operationally meaningful.

\begin{remark}[No numerical overclaim]
These are scaling consequences. Numerical predictions require sector-specific computation of
the relevant operator \(A\), projector \(P\), spectral window \(\chi\), damping gap
\(\lambda_\ast\), admissibility tolerance, and detector coupling data.
\end{remark}

\section{Conclusion}

Wave--particle duality is not a primitive paradox in Modal Triplet Theory. It is the downstream
appearance of one coherent-sector excitation under two different projections.

The particle-like aspect is the local shadow:
\[
  \Kadm(x,y)\to\delta(x-y).
\]
The wave-like aspect is the spectral shadow:
\[
  \psi_{\coh}(x,t)
  =
  \sum_{n\in\coh}
  c_ne^{-i\omega_nt}\phi_n(x).
\]
Measurement is finite survivor-basin selection:
\[
  \rho_{ab}\mapsto D_{ab}\rho_{ab},
\]
implemented, when closed, by a completely positive trace-preserving Schur channel and a
normalized detector POVM.

The double-slit experiment then becomes transparent. When both branches remain coherently
admissible, interference appears. When detector coupling separates the branches into
distinguishable survivor basins, interference is damped and localized events remain.

Thus the quantum object is not forced to be either a wave or a particle. In MTT, both are
downstream shadows of one finite coherent excitation, revealed differently by spectral
evolution, local probing, and admissible selection.

\[
  \boxed{\text{Wave--particle duality is projection duality.}}
\]

\begin{thebibliography}{20}

\bibitem{ReedSimon}
M. Reed and B. Simon,
\emph{Methods of Modern Mathematical Physics I: Functional Analysis},
Academic Press.

\bibitem{Davies}
E. B. Davies,
\emph{Heat Kernels and Spectral Theory},
Cambridge University Press.

\bibitem{GlimmJaffe}
J. Glimm and A. Jaffe,
\emph{Quantum Physics: A Functional Integral Point of View},
Springer.

\bibitem{Zurek}
W. H. Zurek,
Decoherence, einselection, and the quantum origins of the classical,
\emph{Reviews of Modern Physics} 75, 715--775.

\bibitem{Schlosshauer}
M. Schlosshauer,
\emph{Decoherence and the Quantum-to-Classical Transition},
Springer.

\bibitem{Bohr}
N. Bohr,
The quantum postulate and the recent development of atomic theory,
\emph{Nature} 121, 580--590.

\bibitem{Bohm}
D. Bohm,
A suggested interpretation of the quantum theory in terms of hidden variables,
\emph{Physical Review} 85, 166--193.

\bibitem{FeynmanHibbs}
R. P. Feynman and A. R. Hibbs,
\emph{Quantum Mechanics and Path Integrals},
McGraw--Hill.

\bibitem{MTTDelta}
P. Nero,
\emph{Dirac Delta Functions as Singular Shadows of Admissible Projection},
MTT Delta--Projection sequence.

\bibitem{MTTKernel}
P. Nero,
\emph{Canonical Coherent Kernels from MTT Fixed-Point Data},
MTT Delta--Projection sequence.

\bibitem{MTTDamping}
P. Nero,
\emph{Deriving the MTT Coherence Scale from Fixed-Point Damping},
MTT Delta--Projection sequence.


\bibitem{Schoenberg}
I. J. Schoenberg,
Metric spaces and positive definite functions,
\emph{Transactions of the American Mathematical Society} 44, 522--536.

\bibitem{Englert}
B.-G. Englert,
Fringe visibility and which-way information: An inequality,
\emph{Physical Review Letters} 77, 2154--2157.

\bibitem{HornJohnson}
R. A. Horn and C. R. Johnson,
\emph{Matrix Analysis},
Cambridge University Press.

\bibitem{Paulsen}
V. Paulsen,
\emph{Completely Bounded Maps and Operator Algebras},
Cambridge University Press.

\bibitem{NielsenChuang}
M. A. Nielsen and I. L. Chuang,
\emph{Quantum Computation and Quantum Information},
Cambridge University Press.

\bibitem{Polchinski}
J. Polchinski,
\emph{String Theory, Volumes 1 and 2},
Cambridge University Press.

\bibitem{MTTBornBridge}
P. Nero,
\emph{Inflationary Measures and the Born Rule as a Single Shadow--Bridge Problem},
MTT Shadow--Bridge sequence.

\bibitem{MTTDisturbance}
P. Nero,
\emph{Measurement as Disturbance and Stabilization in Modal Triplet Theory},
MTT Measurement sequence.

\bibitem{MTTMeasurement}
P. Nero,
\emph{Measurement Effects as Finite Survivor-Basin Kernels},
MTT Delta--Projection sequence.

\end{thebibliography}

\end{document}
